\section{Projective resolutions}
\label{section-projective-resolutions}

\noindent
This section is dual to
Section \ref{section-injective-resolutions}.
We give definitions and state results, but we do not reprove the lemmas.

\begin{definition}
\label{definition-projective-resolution}
Let $\mathcal{A}$ be an abelian category.
Let $A \in \Ob(\mathcal{A})$.
A {\it projective resolution of $A$} is a complex
$P^\bullet$ together with a map $P^0 \to A$ such
that:
\begin{enumerate}
\item We have $P^n = 0$ for $n > 0$.
\item Each $P^n$ is a projective object of $\mathcal{A}$.
\item The map $P^0 \to A$ induces an isomorphism $\Coker(d^{-1}) \to A$.
\item We have $H^i(P^\bullet) = 0$ for $i < 0$.
\end{enumerate}
Hence $P^\bullet \to A[0]$ is a quasi-isomorphism.
In other words the complex
$$
\ldots \to P^{-1} \to P^0 \to A \to 0 \to \ldots
$$
is acyclic. Let $K^\bullet$ be a complex in $\mathcal{A}$.
A {\it projective resolution of $K^\bullet$} is a complex
$P^\bullet$ together with a map $\alpha : P^\bullet \to K^\bullet$
of complexes such that
\begin{enumerate}
\item We have $P^n = 0$ for $n \gg 0$, i.e., $P^\bullet$ is bounded above.
\item Each $P^n$ is a projective object of $\mathcal{A}$.
\item The map $\alpha : P^\bullet \to K^\bullet$ is a
quasi-isomorphism.
\end{enumerate}
\end{definition}

\begin{lemma}
\label{lemma-cohomology-bounded-above}
Let $\mathcal{A}$ be an abelian category.
Let $K^\bullet$ be a complex of $\mathcal{A}$.
\begin{enumerate}
\item If $K^\bullet$ has a projective resolution then
$H^n(K^\bullet) = 0$ for $n \gg 0$.
\item If $H^n(K^\bullet) = 0$ for $n \gg 0$ then there
exists a quasi-isomorphism $L^\bullet \to K^\bullet$
with $L^\bullet$ bounded above.
\end{enumerate}
\end{lemma}

\begin{proof}
Dual to
Lemma \ref{lemma-cohomology-bounded-below}.
\end{proof}

\begin{lemma}
\label{lemma-projective-resolutions-exist}
Let $\mathcal{A}$ be an abelian category.
Assume $\mathcal{A}$ has enough projectives.
\begin{enumerate}
\item Any object of $\mathcal{A}$ has a projective resolution.
\item If $H^n(K^\bullet) = 0$ for all $n \gg 0$ then
$K^\bullet$ has a projective resolution.
\item If $K^\bullet$ is a complex with $K^n = 0$ for $n > a$, then
there exists a projective resolution $\alpha : P^\bullet \to K^\bullet$
with $P^n = 0$ for $n > a$ such that each $\alpha^n : P^n \to K^n$ is
surjective.
\end{enumerate}
\end{lemma}

\begin{proof}
Dual to
Lemma \ref{lemma-injective-resolutions-exist}.
\end{proof}

\begin{lemma}
\label{lemma-projective-into-acyclic-is-zero}
Let $\mathcal{A}$ be an abelian category.
Let $K^\bullet$ be an acyclic complex.
Let $P^\bullet$ be bounded above and consisting of projective objects.
Any morphism $P^\bullet \to K^\bullet$ is homotopic to zero.
\end{lemma}

\begin{proof}
Dual to
Lemma \ref{lemma-acyclic-is-zero}.
\end{proof}

\begin{remark}
\label{remark-easier-projective}
Let $\mathcal{A}$ be an abelian category.
Suppose that $\alpha : K^\bullet \to L^\bullet$ is a quasi-isomorphism
of complexes. Let $P^\bullet$ be a bounded above complex of projectives.
Then
$$
\Hom_{K(\mathcal{A})}(P^\bullet, K^\bullet)
\longrightarrow
\Hom_{K(\mathcal{A})}(P^\bullet, L^\bullet)
$$
is an isomorphism. This is dual to
Remark \ref{remark-easier-proofs}.
\end{remark}

\begin{lemma}
\label{lemma-morphisms-lift-projective}
Let $\mathcal{A}$ be an abelian category.
Consider a solid diagram
$$
\xymatrix{
K^\bullet & L^\bullet \ar[l]^\alpha \\
P^\bullet \ar[u] \ar@{-->}[ru]_\beta
}
$$
where $P^\bullet$ is bounded above and consists of projective
objects, and $\alpha$ is a quasi-isomorphism.
\begin{enumerate}
\item There exists a map of complexes $\beta$ making the diagram
commute up to homotopy.
\item If $\alpha$ is surjective in every degree
then we can find a $\beta$ which makes the diagram commute.
\end{enumerate}
\end{lemma}

\begin{proof}
Dual to
Lemma \ref{lemma-morphisms-lift}.
\end{proof}

\begin{lemma}
\label{lemma-morphisms-equal-up-to-homotopy-projective}
Let $\mathcal{A}$ be an abelian category. Consider a solid diagram
$$
\xymatrix{
K^\bullet & L^\bullet \ar[l]^\alpha \\
P^\bullet \ar[u] \ar@{-->}[ru]_{\beta_i}
}
$$
where $P^\bullet$ is bounded above and consists of projective
objects, and $\alpha$ is a quasi-isomorphism.
Any two morphisms $\beta_1, \beta_2$ making the diagram commute
up to homotopy are homotopic.
\end{lemma}

\begin{proof}
Dual to
Lemma \ref{lemma-morphisms-equal-up-to-homotopy}.
\end{proof}

\begin{lemma}
\label{lemma-morphisms-from-projective-complex}
Let $\mathcal{A}$ be an abelian category.
Let $P^\bullet$ be bounded above complex consisting of projective
objects. Let $L^\bullet \in K(\mathcal{A})$. Then
$$
\Mor_{K(\mathcal{A})}(P^\bullet, L^\bullet)
=
\Mor_{D(\mathcal{A})}(P^\bullet, L^\bullet).
$$
\end{lemma}

\begin{proof}
Dual to
Lemma \ref{lemma-morphisms-into-injective-complex}.
\end{proof}

\begin{lemma}
\label{lemma-projective-resolution-ses}
Let $\mathcal{A}$ be an abelian category.
Assume $\mathcal{A}$ has enough projectives.
For any short exact sequence
$0 \to A^\bullet \to B^\bullet \to C^\bullet \to 0$
of $\text{Comp}^{+}(\mathcal{A})$ there exists a
commutative diagram in $\text{Comp}^{+}(\mathcal{A})$
$$
\xymatrix{
0 \ar[r] &
P_1^\bullet \ar[r] \ar[d] &
P_2^\bullet \ar[r] \ar[d] &
P_3^\bullet \ar[r] \ar[d] &
0 \\
0 \ar[r] &
A^\bullet \ar[r] &
B^\bullet \ar[r] &
C^\bullet \ar[r] &
0
}
$$
where the vertical arrows are projective resolutions and
the rows are short exact sequences of complexes.
In fact, given any projective resolution $P^\bullet \to C^\bullet$
we may assume $P_3^\bullet = P^\bullet$.
\end{lemma}

\begin{proof}
Dual to
Lemma \ref{lemma-injective-resolution-ses}.
\end{proof}

\begin{lemma}
\label{lemma-precise-vanishing}
Let $\mathcal{A}$ be an abelian category.
Let $P^\bullet$, $K^\bullet$ be complexes.
Let $n \in \mathbf{Z}$. Assume that
\begin{enumerate}
\item $P^\bullet$ is a bounded complex consisting of projective
objects,
\item $P^i = 0$ for $i < n$, and
\item $H^i(K^\bullet) = 0$ for $i \geq n$.
\end{enumerate}
Then
$\Hom_{K(\mathcal{A})}(P^\bullet, K^\bullet) =
\Hom_{D(\mathcal{A})}(P^\bullet, K^\bullet) = 0$.
\end{lemma}

\begin{proof}
The first equality follows from
Lemma \ref{lemma-morphisms-from-projective-complex}.
Note that there is a distinguished triangle
$$
(\tau_{\leq n - 1}K^\bullet, K^\bullet, \tau_{\geq n}K^\bullet, f, g, h)
$$
by Remark \ref{remark-truncation-distinguished-triangle}. Hence, by
Lemma \ref{lemma-representable-homological}
it suffices to prove
$\Hom_{K(\mathcal{A})}(P^\bullet, \tau_{\leq n - 1}K^\bullet) = 0$ and
$\Hom_{K(\mathcal{A})}(P^\bullet, \tau_{\geq n} K^\bullet) = 0$.
The first vanishing is trivial and the second is
Lemma \ref{lemma-projective-into-acyclic-is-zero}.
\end{proof}

\begin{lemma}
\label{lemma-lift-map}
Let $\mathcal{A}$ be an abelian category.
Let $\beta : P^\bullet \to L^\bullet$ and
$\alpha : E^\bullet \to L^\bullet$ be
maps of complexes. Let $n \in \mathbf{Z}$. Assume
\begin{enumerate}
\item $P^\bullet$ is a bounded complex of projectives and
$P^i = 0$ for $i < n$,
\item $H^i(\alpha)$ is an isomorphism for $i > n$ and surjective
for $i = n$.
\end{enumerate}
Then there exists a map of complexes $\gamma : P^\bullet \to E^\bullet$
such that $\alpha \circ \gamma$ and $\beta$ are homotopic.
\end{lemma}

\begin{proof}
Consider the cone $C^\bullet = C(\alpha)^\bullet$ with map
$i : L^\bullet \to C^\bullet$.
Note that $i \circ \beta$ is zero by
Lemma \ref{lemma-precise-vanishing}.
Hence we can lift $\beta$ to $E^\bullet$ by
Lemma \ref{lemma-representable-homological}.
\end{proof}












